\documentclass[10pt,a4paper]{article}
\usepackage[utf8]{inputenc}
\usepackage[T1]{fontenc}
\usepackage[spanish,es-nodecimaldot,es-noquoting]{babel}
\usepackage{lmodern}
\usepackage{amsmath,amssymb,amsfonts,mathtools}
\usepackage[a4paper, top=1.35cm, bottom=1.45cm, left=1.45cm, right=1.45cm, headheight=14pt]{geometry}
\usepackage{array,tabularx,booktabs,multicol,multirow}
\usepackage[table]{xcolor}
\usepackage[most]{tcolorbox}
\usepackage{titlesec}
\usepackage{fancyhdr}
\usepackage{lastpage}
\usepackage[hidelinks, bookmarks=true, pdftitle={Tabla Maestra de Derivadas}, pdfsubject={Reglas de derivacion, tabla simple y compuesta, casos especiales y ejemplos}]{hyperref}

% Paleta sobria consistente con los apuntes de integrales
\definecolor{mainblue}{RGB}{24, 76, 120}
\definecolor{subblue}{RGB}{41, 105, 160}
\definecolor{softblue}{RGB}{244, 248, 253}
\definecolor{borderblue}{RGB}{185, 204, 224}
\definecolor{tipgreen}{RGB}{36, 115, 72}
\definecolor{tipbg}{RGB}{244, 251, 246}
\definecolor{alertred}{RGB}{168, 44, 40}
\definecolor{alertbg}{RGB}{254, 246, 245}
\definecolor{slatetext}{RGB}{55, 65, 81}

\titleformat{\section}{\large\bfseries\color{mainblue}}{\thesection.}{0.5em}{}[\vspace{-3pt}\titlerule]
\titlespacing*{\section}{0pt}{6pt}{4pt}

\setlength{\parindent}{0pt}
\setlength{\parskip}{2.5pt}

\pagestyle{fancy}
\fancyhf{}
\fancyfoot[L]{\footnotesize\color{slatetext}Matemáticas --- Tabla Maestra de Derivadas}
\fancyfoot[R]{\footnotesize\color{slatetext}Página \thepage\ de \pageref{LastPage}}
\renewcommand{\headrulewidth}{0pt}
\renewcommand{\footrulewidth}{0.4pt}

\newtcolorbox{formulabox}[1]{
  unbreakable,
  enhanced,
  colback=softblue,
  colframe=mainblue,
  arc=1.3mm,
  boxrule=0.75pt,
  left=3.2mm, right=3.2mm, top=1.8mm, bottom=1.8mm,
  title={\textbf{#1}},
  coltitle=white,
  colbacktitle=mainblue,
  fonttitle=\small\bfseries,
  before skip=3pt, after skip=4pt
}

\newtcolorbox{trucobox}[1]{
  unbreakable,
  enhanced,
  colback=tipbg,
  colframe=tipgreen,
  arc=1.3mm,
  boxrule=0.75pt,
  left=3.2mm, right=3.2mm, top=1.8mm, bottom=1.8mm,
  title={\textbf{#1}},
  coltitle=white,
  colbacktitle=tipgreen,
  fonttitle=\small\bfseries,
  before skip=3pt, after skip=4pt
}

\newtcolorbox{trampabox}[1]{
  unbreakable,
  enhanced,
  colback=alertbg,
  colframe=alertred,
  arc=1.3mm,
  boxrule=0.75pt,
  left=3.2mm, right=3.2mm, top=1.8mm, bottom=1.8mm,
  title={\textbf{#1}},
  coltitle=white,
  colbacktitle=alertred,
  fonttitle=\small\bfseries,
  before skip=3pt, after skip=4pt
}

\begin{document}

% ==============================================================================
% PAGINA 1: HOJA-CHULETA MAESTRA DE DERIVADAS (IMPRIMIBLE EN 1 CARA)
% ==============================================================================
\begin{center}
  {\LARGE \textbf{\color{mainblue} Tabla Maestra de Derivadas y Reglas de Derivación}}\\[1.5mm]
  {\small \color{slatetext} Formulación elemental, Regla de la Cadena ($u = u(x)$), casos especiales y ejemplos resueltos}
\end{center}
\vspace{1mm}
\hrule
\vspace{3mm}

% --- REGLAS OPERATIVAS ---
\begin{formulabox}{1. Reglas Operativas (Álgebra de Derivadas y Regla de la Cadena)}
\small
\renewcommand{\arraystretch}{1.28}
\begin{tabularx}{\linewidth}{@{}l X l X@{}}
\textbf{Constante por función:} & $(k\cdot f)' = k\cdot f'$ & \textbf{Suma y resta:} & $(f \pm g)' = f' \pm g'$ \\
\textbf{Producto:} & $(f \cdot g)' = f'\cdot g + f\cdot g'$ & \textbf{Cociente:} & $\displaystyle \left(\frac{f}{g}\right)' = \frac{f'\cdot g - f\cdot g'}{g^2}$ \\
\textbf{Producto triple:} & $(fgh)' = f'gh + fg'h + fgh'$ & \textbf{Regla de la Cadena:} & $\big[f(u(x))\big]' = f'(u(x))\cdot u'(x)$ \\
\end{tabularx}
\end{formulabox}

% --- TABLA MAESTRA COMPLETA ---
\begin{center}
\renewcommand{\arraystretch}{1.34}
\footnotesize
\begin{tabularx}{\linewidth}{@{}l >{\raggedright\arraybackslash}X >{\raggedright\arraybackslash}X l@{}}
\toprule
\textbf{Familia / Función} & \textbf{Forma Simple: $y = f(x) \;\longrightarrow\; y'$} & \textbf{Forma Compuesta: $y = f(u) \;\longrightarrow\; y'$} & \textbf{Condición / Nota} \\
\midrule
Constante & $y = k \;\longrightarrow\; y' = 0$ & $y = k \;\longrightarrow\; y' = 0$ & $k \in \mathbb{R}$ \\
Identidad / Lineal & $y = x \;\longrightarrow\; y' = 1$ & $y = ax + b \;\longrightarrow\; y' = a$ & Pendiente constante \\
Potencia & $y = x^n \;\longrightarrow\; y' = n\,x^{n-1}$ & $y = u^n \;\longrightarrow\; y' = n\,u^{n-1}\cdot u'$ & Válida para todo $n \in \mathbb{R}$ \\
Inversa ($n = -1$) & $\displaystyle y = \frac{1}{x} \;\longrightarrow\; y' = -\frac{1}{x^2}$ & $\displaystyle y = \frac{1}{u} \;\longrightarrow\; y' = -\frac{u'}{u^2}$ & Más rápida que cociente \\
Raíz cuadrada & $\displaystyle y = \sqrt{x} \;\longrightarrow\; y' = \frac{1}{2\sqrt{x}}$ & $\displaystyle y = \sqrt{u} \;\longrightarrow\; y' = \frac{u'}{2\sqrt{u}}$ & $u > 0$ \\
Raíz $n$-ésima & $\displaystyle y = \sqrt[n]{x} \;\longrightarrow\; y' = \frac{1}{n\sqrt[n]{x^{n-1}}}$ & $\displaystyle y = \sqrt[n]{u} \;\longrightarrow\; y' = \frac{u'}{n\sqrt[n]{u^{n-1}}}$ & O pasar a $u^{1/n}$ \\
\midrule
Exponencial base $e$ & $y = e^x \;\longrightarrow\; y' = e^x$ & $y = e^u \;\longrightarrow\; y' = e^u \cdot u'$ & $\ln e = 1$ \\
\rowcolor{softblue}
\textbf{Exponencial base $a$} & $\mathbf{y = a^x \;\longrightarrow\; y' = a^x \cdot \ln a}$ & $\mathbf{y = a^u \;\longrightarrow\; y' = a^u \cdot \ln a \cdot u'}$ & $a > 0,\; a \neq 1$ \\
Logaritmo neperiano & $\displaystyle y = \ln x \;\longrightarrow\; y' = \frac{1}{x}$ & $\displaystyle y = \ln u \;\longrightarrow\; y' = \frac{u'}{u}$ & $u > 0$ (también $\ln|u|$) \\
Logaritmo base $a$ & $\displaystyle y = \log_a x \;\longrightarrow\; y' = \frac{1}{x\,\ln a}$ & $\displaystyle y = \log_a u \;\longrightarrow\; y' = \frac{u'}{u\,\ln a}$ & $a > 0,\; a \neq 1,\; u > 0$ \\
\midrule
Seno & $y = \sin x \;\longrightarrow\; y' = \cos x$ & $y = \sin u \;\longrightarrow\; y' = \cos(u)\cdot u'$ & Ángulo en radianes \\
Coseno & $y = \cos x \;\longrightarrow\; y' = -\sin x$ & $y = \cos u \;\longrightarrow\; y' = -\sin(u)\cdot u'$ & Cuidado con el signo $-$ \\
Tangente & $\displaystyle y = \tan x \;\longrightarrow\; y' = \frac{1}{\cos^2 x} = 1+\tan^2 x$ & $\displaystyle y = \tan u \;\longrightarrow\; y' = \frac{u'}{\cos^2 u} = (1+\tan^2 u)\,u'$ & $\sec^2 u$, con $\cos u \neq 0$ \\
Cotangente & $\displaystyle y = \cot x \;\longrightarrow\; y' = -\frac{1}{\sin^2 x} = -(1+\cot^2 x)$ & $\displaystyle y = \cot u \;\longrightarrow\; y' = -\frac{u'}{\sin^2 u} = -(1+\cot^2 u)\,u'$ & $-\csc^2 u$, con $\sin u \neq 0$ \\
Secante & $y = \sec x \;\longrightarrow\; y' = \sec x\,\tan x$ & $y = \sec u \;\longrightarrow\; y' = \sec u\,\tan u \cdot u'$ & $\sec u = \frac{1}{\cos u}$ \\
Cosecante & $y = \csc x \;\longrightarrow\; y' = -\csc x\,\cot x$ & $y = \csc u \;\longrightarrow\; y' = -\csc u\,\cot u \cdot u'$ & $\csc u = \frac{1}{\sin u}$ \\
\midrule
Arcoseno & $\displaystyle y = \arcsin x \;\longrightarrow\; y' = \frac{1}{\sqrt{1 - x^2}}$ & $\displaystyle y = \arcsin u \;\longrightarrow\; y' = \frac{u'}{\sqrt{1 - u^2}}$ & $|u| < 1$ \\
Arcocoseno & $\displaystyle y = \arccos x \;\longrightarrow\; y' = -\frac{1}{\sqrt{1 - x^2}}$ & $\displaystyle y = \arccos u \;\longrightarrow\; y' = -\frac{u'}{\sqrt{1 - u^2}}$ & $|u| < 1$ \\
Arcotangente & $\displaystyle y = \arctan x \;\longrightarrow\; y' = \frac{1}{1 + x^2}$ & $\displaystyle y = \arctan u \;\longrightarrow\; y' = \frac{u'}{1 + u^2}$ & Todo $u \in \mathbb{R}$ \\
Arcocotangente & $\displaystyle y = \operatorname{arccot} x \;\longrightarrow\; y' = -\frac{1}{1 + x^2}$ & $\displaystyle y = \operatorname{arccot} u \;\longrightarrow\; y' = -\frac{u'}{1 + u^2}$ & Todo $u \in \mathbb{R}$ \\
\midrule
Hiperbólicas (extra) & $\sinh x \to \cosh x$; \; $\cosh x \to \sinh x$ & $\sinh u \to \cosh(u)\,u'$; \; $\cosh u \to \sinh(u)\,u'$ & $\tanh u \to \frac{u'}{\cosh^2 u}$ \\
\bottomrule
\end{tabularx}
\end{center}

\begin{trucobox}{Por qué $(a^x)' = a^x \ln a$ y $(\log_a x)' = \frac{1}{x \ln a}$ (para no olvidarlo nunca)}
Escribe cualquier base $a > 0$ en base $e$ usando $a = e^{\ln a}$:
\begin{itemize}\itemsep 1pt
  \item \textbf{Exponencial:} $a^x = e^{x\ln a} \;\implies\; (a^x)' = e^{x\ln a}\cdot (\ln a) = \mathbf{a^x \ln a}$. \textit{(En $e^x$ desaparece porque $\ln e = 1$).}
  \item \textbf{Logaritmo:} $\log_a x = \frac{\ln x}{\ln a} \;\implies\; (\log_a x)' = \frac{1}{\ln a}\cdot\frac{1}{x} = \mathbf{\frac{1}{x\ln a}}$. \textit{(En $\ln x$ desaparece porque $\ln e = 1$).}
\end{itemize}
\end{trucobox}

\newpage

% ==============================================================================
% PAGINA 2: CASOS ESPECIALES, TRUCOS Y EJEMPLOS RESUELTOS
% ==============================================================================
\section*{Casos Especiales, Atajos Previos y Ejemplos Resueltos}

\begin{formulabox}{Los 3 Casos Potenciales/Exponenciales: ¡No los Confundas Jamás!}
\renewcommand{\arraystretch}{1.35}
\small
\begin{tabularx}{\linewidth}{@{}l l l X@{}}
\toprule
\textbf{Tipo} & \textbf{Ejemplo} & \textbf{Regla} & \textbf{Resultado} \\
\midrule
\textbf{1. Base variable, exp. constante} & $y = x^5$ & Regla de la potencia & $y' = 5x^4$ \\
\textbf{2. Base constante, exp. variable} & $y = 5^x$ & Regla exponencial & $y' = 5^x \ln 5$ \\
\textbf{3. Base variable Y exp. variable} & $y = f(x)^{g(x)}$ & \textbf{Derivación logarítmica} & $\displaystyle y' = f^g \cdot \left(g'\ln f + g\cdot\frac{f'}{f}\right)$ \\
\bottomrule
\end{tabularx}
\end{formulabox}

\subsection*{Derivación Logarítmica (paso a paso para $y = f(x)^{g(x)}$)}
Cuando tanto la base como el exponente contienen $x$, \textbf{NO} valen ni la regla de la potencia ni la exponencial. Aplica logaritmo neperiano a ambos lados para bajar el exponente:
\begin{enumerate}\itemsep 1.5pt
  \item Toma $\ln$ en ambos lados: $\ln y = \ln\big(f(x)^{g(x)}\big) = g(x)\cdot\ln f(x)$.
  \item Deriva ambos miembros respecto a $x$ (en la izquierda queda $\frac{y'}{y}$ por regla de la cadena):
  $$\frac{y'}{y} = g'(x)\ln f(x) + g(x)\cdot\frac{f'(x)}{f(x)}$$
  \item Despeja $y'$ multiplicando todo por $y = f(x)^{g(x)}$.
\end{enumerate}

\begin{trucobox}{3 Atajos que Ahorran el 70\% de las Cuentas Antes de Derivar}
\begin{itemize}\itemsep 1.5pt
  \item \textbf{Propiedades de logaritmos antes de derivar:} si tienes $\ln\!\left(\frac{A\cdot B}{C^3}\right)$, sepáralo en $\ln A + \ln B - 3\ln C$ y deriva términos sencillísimos en lugar de hacer una cadena gigante.
  \item \textbf{Constante arriba o abajo:} en $\frac{5}{x^3}$, escribe $5x^{-3}$ y deriva como potencia ($-15x^{-4}$); en $\frac{x^2+1}{4}$, saca $\frac{1}{4}(x^2+1)$ y deriva ($\frac{1}{4}\cdot 2x = \frac{x}{2}$). \textbf{No uses la regla del cociente si arriba o abajo solo hay un número.}
  \item \textbf{Valor absoluto:} para $y = |u(x)|$ (con $u(x) \neq 0$), su derivada es $\displaystyle y' = \frac{u(x)}{|u(x)|}\cdot u'(x)$. En los puntos donde $u(x) = 0$ suele haber un pico (no derivable).
\end{itemize}
\end{trucobox}

\subsection*{10 Ejemplos Resueltos de Referencia Rápida}

\begin{center}
\renewcommand{\arraystretch}{1.48}
\footnotesize
\begin{tabularx}{\linewidth}{@{}c l >{\raggedright\arraybackslash}X >{\raggedright\arraybackslash}X@{}}
\toprule
\textbf{\#} & \textbf{Función $y = f(x)$} & \textbf{Identificación y aplicación de la regla} & \textbf{Derivada final $y'$} \\
\midrule
\textbf{1} & $y = 3^x + e^{4x-1}$ & Exponencial base $3$ ($a^x \ln a$) + cadena en $e^u$ ($u'=4$) & $y' = 3^x \ln 3 + 4e^{4x-1}$ \\
\textbf{2} & $y = 2^{x^2 + 3x}$ & $a^u$ con $a=2$, $u = x^2+3x \implies u'=2x+3$ & $y' = (2x+3)\cdot 2^{x^2+3x}\ln 2$ \\
\textbf{3} & $y = \log_5(3x^2 - 1)$ & $\log_a u \to \frac{u'}{u\ln a}$ con $u = 3x^2-1$, $u' = 6x$ & $\displaystyle y' = \frac{6x}{(3x^2-1)\ln 5}$ \\
\textbf{4} & $y = \frac{4}{x^3} - 2\sqrt{x} + \sqrt[3]{x^2}$ & Pasar a potencias: $4x^{-3} - 2x^{1/2} + x^{2/3}$ & $\displaystyle y' = -\frac{12}{x^4} - \frac{1}{\sqrt{x}} + \frac{2}{3\sqrt[3]{x}}$ \\
\textbf{5} & $y = x^2 e^{-3x}$ & Producto: $(x^2)'e^{-3x} + x^2(e^{-3x})' = 2xe^{-3x} - 3x^2 e^{-3x}$ & $y' = (2x - 3x^2)e^{-3x}$ \\
\textbf{6} & $y = \frac{\ln x}{x^2}$ & Cociente: $\frac{(1/x)\cdot x^2 - (\ln x)(2x)}{x^4} = \frac{x(1 - 2\ln x)}{x^4}$ & $\displaystyle y' = \frac{1 - 2\ln x}{x^3}$ \\
\textbf{7} & $y = \sin^3(2x)$ & Cadena en 3 capas: $[\sin(2x)]^3 \to 3\sin^2(2x)\cdot\cos(2x)\cdot 2$ & $y' = 6\sin^2(2x)\cos(2x)$ \\
\textbf{8} & $y = \arctan(e^x)$ & $\arctan u \to \frac{u'}{1+u^2}$ con $u = e^x \implies u' = e^x$ & $\displaystyle y' = \frac{e^x}{1 + e^{2x}}$ \\
\textbf{9} & $y = \ln\!\left(\frac{x^3}{\sqrt{x+1}}\right)$ & Atajo previo: $3\ln x - \frac{1}{2}\ln(x+1)$ y derivar directo & $\displaystyle y' = \frac{3}{x} - \frac{1}{2(x+1)}$ \\
\textbf{10} & $y = x^x \quad (x > 0)$ & Tomar $\ln$: $\ln y = x\ln x \implies \frac{y'}{y} = \ln x + 1$ & $y' = x^x(\ln x + 1)$ \\
\bottomrule
\end{tabularx}
\end{center}

\begin{trampabox}{Los 6 Errores Típicos al Derivar}
\footnotesize
\begin{multicols}{2}
\begin{enumerate}\itemsep 1.5pt
  \item Hacer $(a^x)' = x\,a^{x-1}$ en lugar de $a^x \ln a$.
  \item Olvidar el signo menos en $(\cos u)' = -\sin u \cdot u'$ y $(\arccos u)'$.
  \item Olvidar multiplicar por $u'$ al final de la Regla de la Cadena.
  \item Derivar $(u\cdot v)' = u'\cdot v'$ (¡falso! es $u'v + uv'$).
  \item Escribir $(\ln u)' = \frac{1}{u}$ sin el $u'$ en el numerador.
  \item Derivar $x^x$ sin aplicar derivación logarítmica.
\end{enumerate}
\end{multicols}
\end{trampabox}

\end{document}
